\setcounter{section}{0}

\Needspace{5\baselineskip}
\hypertarget{ux5de5ux5177ux8c03ux7528-1}{%
\section{工具调用}\label{ux5de5ux5177ux8c03ux7528-1}}

为了让大语言模型成为从单纯的``大脑''变为能实时感知并执行任务的智能体，我们需要让它具备调用工具的能力，具体包括以下方面。

\Needspace{5\baselineskip}
\hypertarget{ux4e00ux5de5ux5177ux8c03ux7528ux7684ux57faux672cux6d41ux7a0bux4e0eux7b56ux7565}{%
\subsection{一、工具调用的基本流程与策略}\label{ux4e00ux5de5ux5177ux8c03ux7528ux7684ux57faux672cux6d41ux7a0bux4e0eux7b56ux7565}}

\begin{enumerate}
\def\labelenumi{\arabic{enumi}.}
\item
  基本流程：判断是否需要使用工具---选择合适的工具---调用工具---解析输出结果。
\item
  策略
\end{enumerate}

（1）提示词驱动：如加入prompt``如果你觉得当前知识不足或需要实时信息，请使用对应工具。如果你可以直接回答，就直接生成答案。''提示词中一般会加入可用的工具定义；考虑到上下文有限，实际上一般只加入``元工具''的定义（如bash命令，模型通过运行命令可以调用其他次级工具），或只放工具列表（模型在需要时自主至文件系统搜索工具）。

\Needspace{8\baselineskip}
（2）ReAct范式（Reason + Act）

\Needspace{5\baselineskip}
前面讲了用提示词驱动大模型使用工具，也讲过推理中的思维链，但这两者是并行的，它们之间并没有很好地协同。ReAct范式则将它们结合起来，比如下面的例子：

Question: 一美元等于多少日元？

Thought: 我不知道当前汇率，应该先查一下。

Action: SEARCH{[}一美元对日元汇率{]}

Observation: 当前汇率为 157.5

Thought: 得到了汇率，现在可以计算。

Action: CALC{[}157.5 × 20{]}

Observation: 结果为 3150

Final Answer: 20美元约等于3150日元。

当然，本质上这还是用提示词实现的。

（3）通过RL训练，让模型自主学会在合适的时候调用工具。

\Needspace{5\baselineskip}
\hypertarget{ux4e8cux5de5ux5177ux7684ux7c7bux522b}{%
\subsection{二、工具的类别}\label{ux4e8cux5de5ux5177ux7684ux7c7bux522b}}

\begin{enumerate}
\def\labelenumi{\arabic{enumi}.}
\item
  信息检索类工具：对应人的``感知神经元''，帮助模型实时获得信息或知识，如搜索引擎、知识数据库、上传文档的解析工具等。
\item
  计算推理类工具：对应人的``中枢神经元''，帮助模型精确运算或推理，避免基于概率的大模型自身推理中可能出现的错误，如计算器、代码解释器等。
\item
  设备控制类工具：对应人的``运动神经元''，帮助模型在实际环境中执行任务，如控制操作系统、应用程序、机器人等的工具。
\end{enumerate}

\Needspace{5\baselineskip}
\hypertarget{ux4e09ux5de5ux5177ux7684ux63a5ux53e3}{%
\subsection{三、工具的接口}\label{ux4e09ux5de5ux5177ux7684ux63a5ux53e3}}

为了让模型理解工具的用途、调用方式等，我们需要设计特定的API接口，以用模型能理解的形式将它们封装起来。工具接口需要有工具名称、用途描述、输入参数，以及对输入参数中有哪些必须填写的说明。

\Needspace{5\baselineskip}
\hypertarget{ux56dbux591aux79cdux5de5ux5177ux7684ux6574ux5408}{%
\subsection{四、多种工具的整合}\label{ux56dbux591aux79cdux5de5ux5177ux7684ux6574ux5408}}

\begin{enumerate}
\def\labelenumi{\arabic{enumi}.}
\item
  工具的选择：可以在提示中提供工具列表，或设置一个关键词触发路由等。
\item
  工具优先级与回退机制：优先调用一些成功率更高或代价更低的工具；如果调用失败，切换至优先级更低的工具；失败若干次，则直接回答，避免陷入死循环。还可以引入反思机制，鼓励模型评估失败原因并改进下一步调用。
\end{enumerate}

\Needspace{5\baselineskip}
\hypertarget{ux4e94ux5de5ux5177ux8fd4ux56deux7ed3ux679cux7684ux538bux7f29}{%
\subsection{五、工具返回结果的压缩}\label{ux4e94ux5de5ux5177ux8fd4ux56deux7ed3ux679cux7684ux538bux7f29}}

工具返回的结果有时候巨大无比，会占满上下文空间。此时可以采取剪枝方式，压缩时间靠前的输出后再输入给LLM。这个压缩是有损过程，但这个过程不重写硬盘中的返回结果文件。

\Needspace{5\baselineskip}
\hypertarget{ux516dux601dux7ef4ux4e0aux4e0bux6587ux7ba1ux7406ux89e3ux51b3ux8c03ux7528ux5de5ux5177ux540eux7684ux63a8ux7406ux65adux94fedeepseek2025}{%
\subsection{六、思维上下文管理：解决调用工具后的推理``断链''（DeepSeek，2025）}\label{ux516dux601dux7ef4ux4e0aux4e0bux6587ux7ba1ux7406ux89e3ux51b3ux8c03ux7528ux5de5ux5177ux540eux7684ux63a8ux7406ux65adux94fedeepseek2025}}

常见的含工具调用的推理模型的工作流往往是：用户问题 → 模型``内部思考''（状态A）→ 执行工具 → 工具结果 + 用户问题 → 模型``内部思考''（状态B）。在进入状态B时，状态A并没有被保存起来。

DeepSeek设计了一套让模型能像人类一样``参考笔记''来完成复杂任务的机制。V3.2的流程是：用户问题 → 模型生成 reasoning\_content\_A（写在``草稿纸''上）→ 执行工具 → 工具结果 + 用户问题 + reasoning\_content\_A（作为上下文输入） → 模型基于之前的``草稿''继续推理，生成 reasoning\_content\_B。

\FloatBarrier
\Needspace{5\baselineskip}
\hypertarget{ux53c2ux8003ux6587ux732e-122}{%
\subsection{参考文献}\label{ux53c2ux8003ux6587ux732e-122}}

\vspace{0.25em}
\begin{itemize}
\tightlist
\item
  温睦宁、林江浩、张伟楠、俞勇. (2025). \href{https://haa.boyuai.com/}{《动手学大模型智能体》}. 人民邮电出版社. ISBN 978-7-115-68638-1.
\item
  Yao, S., Zhao, J., Yu, D., et al.~(2023). \href{https://arxiv.org/abs/2210.03629}{ReAct: Synergizing Reasoning and Acting in Language Models}. ICLR 2023.
\item
  DeepSeek-AI. (2025). \href{https://arxiv.org/abs/2512.02556}{DeepSeek-V3.2: Pushing the Frontier of Open Large Language Models}. arXiv:2512.02556.
\end{itemize}

\clearpage
